\chapter*{Introduction}
\label{chap:introduction} 
 In this book we present a complete
and detailed proof of the

\begin{quote}{\bf Poincar\'e Conjecture:
a closed, smooth, simply connected $3$-manifold is
diffeomorphic}\footnote{Every topological $3$-manifold admits a
differentiable structure and every homeomorphism between smooth
$3$-manifolds can be approximated by a diffeomorphism. Thus,
classification results about topological $3$-manifolds up to
homeomorphism and about smooth $3$-manifolds up to diffeomorphism
are equivalent. In this book `manifold' means `smooth manifold.'}
{\bf to $S^3$.}
\end{quote}
%%% quote environment added by JC

\noindent %%% \noindent addd by JC
This conjecture was formulated by Henri Poincar\'e \cite{Poincare}
in 1904 and has remained open until the recent work of Perelman. The
arguments we give here are a detailed version of those that appear
in Perelman's three preprints~\cite{Perelman1,Perelman2,Perelman3}.
Perelman's arguments rest on a foundation built by Richard Hamilton
with his study of the Ricci flow equation for Riemannian metrics.
Indeed, Hamilton believed that Ricci flow could be used to establish
the Poincar\'e Conjecture and more general topological
classification results in dimension $3$, and laid out a program to
accomplish this. The difficulty was to deal with singularities in
the Ricci flow. Perelman's breakthrough was to understand the
qualitative nature of the singularities sufficiently to allow him to
prove the Poincar\'e Conjecture (and Theorem~\ref{Theorem1} below
which implies the Poincar\'e Conjecture).   For
a detailed history of the Poincar\'e Conjecture, see Milnor's survey
article \cite{Milnorsurvey}.

\bigskip
A  class of examples closely related to the $3$-sphere are {\em the
$3$-dimensional spherical space-forms}, i.e.,
the quotients of $S^3$ by free, linear actions of finite subgroups of the
orthogonal group $O(4)$. There is a generalization of the Poincar\'e
Conjecture, called the {\bf $3$-dimensional spherical space-form
conjecture}, which conjectures
that any closed $3$-manifold with finite fundamental group is diffeomorphic to
a $3$-dimensional spherical space-form. Clearly, a special case of the
$3$-dimensional spherical space-form conjecture is the Poincar\'e Conjecture.

 As indicated in Remark 1.4 of \cite{Perelman3},  the arguments we present
 here not only prove the Poincar\'e Conjecture, they prove the $3$-dimensional space-form conjecture.
 In fact,  the purpose of this book is to prove the following
 more general theorem.

\begin{theorem}[Classification of 3-manifolds with free-product fundamental group]
\label{thm:classification-fundamental-group-free-product}\label{Theorem1} Let $M$ be a closed, connected $3$-manifold and
suppose that the fundamental group of $M$ is a free product of
finite groups and infinite cyclic groups. Then $M$ is diffeomorphic
to a connected sum of spherical space-forms, copies of $S^2\times
S^1$, and copies of the unique (up to diffeomorphism) non-orientable
$2$-sphere bundle over $S^1$.
\end{theorem}

This immediately implies an affirmative resolution of the Poincar\'e
Conjecture and of the $3$-dimensional spherical space-form
conjecture.

\begin{corollary}[Poincaré Conjecture and spherical space-form conjecture]
\label{cor:poincare-and-spherical-space-form}
\uses{thm:classification-fundamental-group-free-product}
(a) A closed, simply connected $3$-manifold is diffeomorphic to
$S^3$. (b) A closed  $3$-manifold with finite fundamental
group is diffeomorphic to a $3$-dimensional spherical space-form.
\end{corollary}

Before launching into a more detailed description of the contents of
this book, one remark on the style of the exposition is in order.
Because of the importance and visibility of the results discussed
here, and because of the number of incorrect claims of proofs of
these results in the past, we felt that it behooved us to work out
and present the arguments in great detail. Our goal was to make the
arguments clear and convincing and also to make them more easily
accessible to a wider audience. As a result, experts may find some
of the points are overly elaborated.



\section{Overview of Perelman's argument}


In dimensions less than or equal to three, any Riemannian metric of constant
Ricci curvature has constant sectional curvature. Classical results in
Riemannian geometry show that the universal cover of a closed manifold of
constant positive curvature is diffeomorphic to the sphere and that the
fundamental group is identified with a finite subgroup of the orthogonal group
acting linearly and freely on the universal cover. Thus, one can approach the
Poincar\'e Conjecture and the more general $3$-dimensional spherical space-form
problem by asking the following question. Making the appropriate fundamental
group assumptions on $3$-manifold $M$, how does one establish the existence of
a metric of constant Ricci curvature on $M$? The essential ingredient in
producing such a metric is the Ricci flow equation
introduced by Richard Hamilton in \cite{Hamilton3MPRC}:
$$\frac{\partial g(t)}{\partial t}=-2\Ric(g(t)),$$
where $\Ric(g(t))$ is the Ricci curvature of the metric $g(t)$.
The fixed points (up to rescaling) of this equation are the
Riemannian metrics of constant Ricci curvature. For a general
introduction to the subject of the Ricci flow see Hamilton's survey
paper \cite{Hamiltonsurvey}, the book by Chow-Knopf
\cite{ChowKnopf}, or the book by Chow, Lu, and Ni \cite{ChowLuNi}.
The Ricci flow equation is a (weakly) parabolic partial differential
flow equation for Riemannian metrics on a smooth manifold. Following
Hamilton, one defines a Ricci flow to be a family
of Riemannian metrics $g(t)$ on a fixed smooth manifold,
parameterized by $t$ in some interval, satisfying this equation. One
considers $t$ as time and studies the equation as an initial value
problem: Beginning with any Riemannian manifold $(M,g_0)$ find a
Ricci flow with $(M,g_0)$ as initial metric; that is to say find a
one-parameter family $(M,g(t))$ of Riemannian manifolds with
$g(0)=g_0$ satisfying the Ricci flow equation. This equation is
valid in all dimensions but we concentrate here on dimension three.
In a sentence, the method of proof is to begin with any Riemannian
metric on the given smooth $3$-manifold and flow it using the Ricci
flow equation to obtain the constant curvature metric for which one
is searching. There are two examples where things work in exactly
this way, both due to Hamilton. (i) If the initial metric has
positive Ricci curvature, Hamilton proved over twenty years ago,
\cite{Hamilton3MPRC}, that under the Ricci flow the manifold shrinks
to a point in finite time, that is to say, there is a finite-time
singularity, and, as we approach the singular time, the diameter of
the manifold tends to zero and the curvature blows up at every
point. Hamilton went on to show that, in this case, rescaling by a
time-dependent function so that the diameter is constant produces a
one-parameter family of metrics converging smoothly to a metric of
constant positive curvature. (ii) At the other extreme, in
\cite{HamiltonNSRF3M} Hamilton showed that if the Ricci flow exists
for all time and if there is an appropriate curvature bound together
with another geometric bound, then as $t\rightarrow\infty$, after
rescaling to have a fixed diameter, the metric converges to a metric
of constant negative curvature.

The results in the general case are much more complicated to
formulate and much more difficult to establish. While Hamilton
established that the Ricci flow equation  has short-term existence
properties, i.e., one can  define $g(t)$ for $t$ in some interval
$[0,T)$ where $T$ depends on the initial metric, it turns out that
if the topology of the manifold is sufficiently complicated, say it
is a non-trivial connected sum, then no matter what the initial
metric is one must encounter finite-time singularities, forced by
the topology. More seriously, even if the manifold has simple
topology, beginning with an arbitrary metric one expects to (and
cannot rule out the possibility that one will) encounter finite-time
singularities in the Ricci flow. These singularities, unlike in the
case of positive Ricci curvature, occur along proper subsets of the
manifold, not the entire manifold.
 Thus, to derive the topological
consequences stated above, it is not sufficient in general to stop
the analysis the first time a singularity arises in the Ricci flow.
One is led to study a more general evolution process
 called {\em Ricci flow with surgery},
 first introduced by Hamilton in the context
 of four-manifolds, \cite{Hamilton4MPIC}.
 This evolution process
 is still parameterized by an interval in time, so that for each $t$
 in the interval of definition there is a compact Riemannian $3$-manifold
 $M_t$. But there is a discrete set of
times at which the manifolds and metrics undergo topological and
metric discontinuities (surgeries). In each of the complementary
intervals to the singular times, the evolution is the usual Ricci
flow, though, because of the surgeries, the topological type of the
manifold $M_t$ changes as $t$ moves from one complementary interval
to the next. From an analytic point of view, the surgeries at the
discontinuity times are introduced in order to `cut away' a
neighborhood of the singularities as they develop and insert by
hand, in place of the `cut away' regions,  geometrically nice
regions. This allows one to continue  the Ricci flow (or more
precisely, restart the Ricci flow with the new metric constructed at
the discontinuity time). Of course, the surgery process also changes
the topology. To be able to say anything useful topologically about
such a process, one needs results about Ricci flow, and one also
needs to control both the topology and the geometry of the surgery
process at the singular times. For example, it is crucial for the
topological applications that we do surgery along $2$-spheres rather
than surfaces of higher genus. Surgery along $2$-spheres produces
the connected sum decomposition, which is well-understood
topologically, while, for example, Dehn surgeries along tori can
completely destroy the topology, changing any $3$-manifold into any
other.

The change in topology turns out to be completely understandable and
amazingly, the surgery processes produce exactly the topological
operations needed to cut the manifold into pieces on which the Ricci
flow can produce the metrics sufficiently controlled so that the
topology can be recognized.

The bulk of this book (Chapters 1-17 and the Appendix) concerns the
establishment of the following long-time existence result for Ricci
flow with surgery.

\begin{theorem}[Long-time existence of Ricci flow with surgery]
\label{thm:ricci-flow-surgery-existence}\label{surgery} Let $(M,g_0)$ be a closed Riemannian $3$-manifold.
Suppose that there is no embedded, locally separating $\Ar P^2$
contained\footnote{I.e., no embedded $\Ar P^2$ in $M$ with trivial
normal bundle. Clearly, all orientable manifolds satisfy this
condition.} in $M$. Then there is a Ricci flow with surgery defined
for all $t\in [0,\infty)$ with initial metric $(M,g_0)$. The set of
discontinuity times for this Ricci flow with surgery is a discrete
subset of $[0,\infty)$. The topological change in the $3$-manifold
as one crosses a surgery time is a connected sum decomposition
together with removal of connected components, each of which is
diffeomorphic to one of $S^2\times S^1$, $\Ar P^3\#\Ar P^3$, the
non-orientable $2$-sphere bundle over $S^1$, or a manifold admitting
a metric of constant positive curvature.
\end{theorem}




While Theorem~\ref{surgery} is central for all applications of Ricci
flow to the topology of three-dimensional manifolds, the argument
for the $3$-manifolds described in Theorem~\ref{Theorem1} is
simplified, and avoids all references to the nature of the flow  as
time goes to infinity, because of the following finite-time
extinction result.

\begin{theorem}[Finite-time extinction of Ricci flow with surgery]
\label{thm:ricci-flow-surgery-finite-extinction}
\uses{thm:ricci-flow-surgery-existence}\label{finiteext}
Let $M$ be a closed $3$-manifold whose fundamental group is a free
product of finite groups and infinite cyclic groups\footnote{In
\cite{Perelman3} Perelman states the result for $3$-manifolds
without prime factors that are acyclic. It is a standard exercise in
$3$-manifold topology to show that Perelman's condition is
equivalent to the group theory hypothesis stated here; see
Corollary~\ref{corequiv}.}. Let $g_0$ be any Riemannian metric on
$M$. Then $M$ admits no locally separating $\Ar P^2$, so that there
is a Ricci flow with surgery defined for all positive time with
$(M,g_0)$ as initial metric as described in Theorem~\ref{surgery}.
This Ricci flow with surgery becomes extinct after some time
$T<\infty$, in the sense that the manifolds $M_t$ are empty for all
$t\ge T$.
\end{theorem}

This result is established in Chapter 18 following the argument
given by Perelman in \cite{Perelman3}, see also
\cite{ColdingMinicozzi}.



We immediately deduce Theorem~\ref{Theorem1} from
Theorems~\ref{surgery} and~\ref{finiteext} as  follows: Let $M$ be a
$3$-manifold satisfying the hypothesis of Theorem~\ref{Theorem1}.
Then there is a finite sequence $M=M_0,M_1,\ldots,M_k=\emptyset$
such that for each $i,\ 1\le i\le k$, $M_i$ is obtained from
$M_{i-1}$ by  a connected sum decomposition or $M_i$ is obtained
from $M_{i-1}$ by removing  a component diffeomorphic to one of
$S^2\times S^1$, $\Ar P^3\#\Ar P^3$, a non-orientable $2$-sphere
bundle over $S^1$, or a $3$-dimensional spherical space-form.
Clearly, it follows by downward induction on $i$ that  each
connected component of $M_i$ is diffeomorphic to a connected sum of
$3$-dimensional spherical space-forms, copies of $S^2\times S^1$,
and copies of the non-orientable $2$-sphere bundle over $S^1$. In
particular, $M=M_0$ has this form. Since $M$ is connected by
hypothesis, this proves the theorem. In fact, this argument proves
the following:

\begin{corollary}[Equivalent conditions for finite-time extinction]
\label{cor:extinction-equivalence}\label{corequiv}
Let $(M_0,g_0)$ a connected Riemannian manifold with no locally
separating $\Ar P^2$. Let $({\mathcal M},G)$ be a Ricci flow with
surgery defined for $0\le t<\infty$ with $(M_0,g_0)$ as initial
manifold. Then the following four conditions are equivalent:
\begin{enumerate} \item[(1)] $({\mathcal M},G)$ becomes extinct after a
finite time, i.e., $M_T=\emptyset$ for all $T$ sufficiently large,
\item[(2)] $M_0$ is diffeomorphic to a connected sum of three-dimensional
spherical space-forms and $S^2$-bundles over $S^1$,\item[(3)] the
fundamental group of $M_0$ is a free product of finite groups and
infinite cyclic groups,
\item[(4)] no prime\footnote{A three-manifold $P$ is prime
if  every separating two-sphere in $P$ bounds a three-ball in $P$.
Equivalently, $P$ is prime if it admits no non-trivial connected sum
decomposition. Every closed three-manifold decomposes as a connected
sum of prime factors with the decomposition being unique up to
diffeomorphism of the factors and the order of the factors.} factor
of $M_0$ is acyclic, i.e., every prime factor of $M_0$ has either
non-trivial $\pi_2$ or non-trivial $\pi_3$.
\end{enumerate}
\end{corollary}

\begin{proof}
\uses{thm:ricci-flow-surgery-existence,thm:ricci-flow-surgery-finite-extinction}
Repeated application of Theorem~\ref{surgery} shows that (1) implies
(2). The implication (2) implies (3) is immediate from van Kampen's
theorem. The fact that (3) implies (1) is Theorem~\ref{finiteext}.
This shows that (1), (2) and (3) are all equivalent. Since
three-dimensional spherical space-forms and $S^2$-bundles over $S^1$
are easily seen to be prime,  (2) implies  (4). Thus, it remains
only to see that (4) implies (3). We consider a manifold $M$
satisfying (4), a prime factor $P$ of $M$, and  universal covering
$\widetilde P$ of $P$. First suppose that $\pi_2(P)=\pi_2(\widetilde
P)$ is trivial. Then, by hypothesis $\pi_3(P)=\pi_3(\widetilde P)$
is non-trivial. By the Hurewicz theorem this means that
$H_3(\widetilde P)$ is non-trivial, and hence that $\widetilde P$ is
a compact, simply connected three-manifold. It follows that
$\pi_1(P)$ is finite. Now suppose that $\pi_2(P)$ is non-trivial.
Then $P$ is not diffeomorphic to $\Ar P^3$. Since $P$ is prime and
contains no locally separating $\Ar P^2$, it follows that $P$
contains no embedded $\Ar P^2$. Then by the sphere theorem there is
an embedded two-sphere in $P$ that is homotopically non-trivial.
Since $P$ is prime, this sphere cannot separate, so cutting $P$ open
along it produces a connected manifold $P_0$ with two boundary
two-spheres. Since $P_0$ is prime, it follows that $P_0$ is
diffeomorphic to $S^2\times I$ and hence $P$ is diffeomorphic to a
two-sphere bundle over the circle.
\end{proof}

\begin{remark}[Remarks on finite-time extinction and its consequences]
\label{rem:extinction-remarks}
\uses{thm:classification-fundamental-group-free-product,cor:extinction-equivalence}
(i) The use of the sphere theorem is unnecessary in the above
argument for what we actually prove is that if every prime factor of
$M$ has non-trivial $\pi_2$ or non-trivial $\pi_3$, then the Ricci
flow with surgery with $(M,g_0)$ as initial metric becomes extinct
after a finite time. In fact, the sphere theorem for closed
three-manifolds follows from the results here.

\noindent (ii)  If the initial manifold is simpler then all the
time-slices are simpler: If $({\mathcal M},G)$ is a Ricci flow with
surgery whose initial manifold is prime, then every time-slice is a
disjoint union of connected components, all but at most one being
diffeomorphic to a three-sphere and if there is one not
diffeomorphic to a three-sphere, then it is diffeomorphic to the
initial manifold. If the initial manifold is a simply connected
manifold $M_0$, then every component of every time-slice $M_T$ must
be simply connected, and thus {\em a posteriori} every time-slice is
a disjoint union of manifolds diffeomorphic to the three-sphere.
Similarly, if the initial manifold has finite fundamental group,
then every connected component of every time-slice is either simply
connected or has the same fundamental group as the initial manifold.

\noindent (iii) The conclusion of this result is a natural
generalization of Hamilton's conclusion in analyzing the Ricci flow
on manifolds of positive Ricci curvature in \cite{Hamilton3MPRC}.
Namely, under appropriate hypotheses, during the evolution process
of Ricci flow with surgery the manifold breaks into components each
of which disappears in finite time. As a component disappears at
some finite time, the metric on that component is well enough
controlled to show that the disappearing component admits a
non-flat, homogeneous Riemannian metric of non-negative sectional
curvature, i.e., a metric locally isometric to either a round $S^3$
or to a product of a round $S^2$ with the usual metric on $ \Ar$.
The existence of such a metric on a component immediately gives the
topological conclusion of Theorem~\ref{Theorem1} for that component,
i.e., that it is diffeomorphic to a $3$-dimensional spherical
space-form, to $S^2\times S^1$ to a non-orientable $2$-sphere bundle
over $S^1$, or to $\Ar P^3\#\Ar P^3$. The biggest difference between
these two results is that Hamilton's hypothesis is geometric
(positive Ricci curvature) whereas Perelman's is homotopy theoretic
(information about the fundamental group).

\noindent (iv) It is also worth pointing out that it follows from
Corollary~\ref{corequiv} that the manifolds that satisfy the four equivalent
conditions in that corollary are exactly the closed, connected, three-manifolds
that admit a Riemannian metric of positive scalar curvature, cf,
\cite{SchoenYau2} and \cite {GromovLawson}.
\end{remark}



One can use Ricci flow in a more general study of three-manifolds
than the one we carry out here. There is a conjecture due to
Thurston, see \cite{Thurstongeom},  known as Thurston's
Geometrization Conjecture or simply as the Geometrization Conjecture for
three-manifolds. It conjectures that every $3$-manifold without
locally separating $\Ar P^2$'s (in particular every orientable
$3$-manifold) is a connected sum of prime $3$-manifolds each of
which admits a decomposition along incompressible\footnote{I.e.,
embedded by a map that is injective on $\pi_1$.} tori into pieces
that admit locally homogeneous geometries of finite volume.  Modulo
questions about cofinite-volume lattices in $SL_2(\Cee)$, proving
this conjecture leads to a complete classification
 of  $3$-manifolds without locally separating $\Ar P^2$'s,
and in particular to a complete classification of all orientable
$3$-manifolds. (See Peter Scott's survey article \cite{Scott}.) By
passing to the orientation double cover and working equivariantly,
these results can be extended to all $3$-manifolds.


 Perelman in
\cite{Perelman2} has stated results which imply a positive
resolution of Thurston's Geometrization conjecture. Perelman's
proposed proof of Thurston's Geometrization Conjecture relies in an
essential way on Theorem~\ref{surgery}, namely the existence of
Ricci flow with surgery for all positive time. But it also involves
a further analysis of the limits of these Ricci flows as time goes
to infinity. This further analysis involves analytic arguments which
are exposed in Sections 6 and 7 of Perelman's second paper
(\cite{Perelman2}), following earlier work of Hamilton
(\cite{HamiltonNSRF3M}) in a simpler case of bounded curvature. They
also involve a result (Theorem 7.4 from \cite{Perelman2}) from the
theory of manifolds with curvature locally bounded below that are
collapsed, related to results of Shioya-Yamaguchi \cite{SY}. The
Shioya-Yamaguchi results in turn rely on an earlier, unpublished
work of Perelman proving the so-called `Stability Theorem.'
Recently, Kapovich, \cite{Kap} has put a preprint on the archive
giving a proof of the stability result. We have been examining
another approach, one suggested by Perelman in \cite{Perelman2},
avoiding the stability theorem, cf, \cite{KleinerLott2} and
\cite{MorganTian2}. It is our view that the collapsing results
needed for the Geometrization Conjecture are in place, but that
before a definitive statement that the Geometrization Conjecture has
been resolved can be made these arguments must be subjected to the
same close scrutiny that the arguments proving the Poincar\'e
Conjecture have received. This process is underway.

In this book we do not attempt to explicate any of the results
beyond Theorem~\ref{surgery} described in the previous paragraph
that are needed for the Geometrization Conjecture. Rather, we
content ourselves with presenting a proof of Theorem~\ref{Theorem1}
above which, as we have indicated, concerns initial Riemannian
manifolds for which the Ricci flow with surgery becomes extinct
after finite time. We are currently preparing a detailed proof,
along the lines suggested by Perelman, of the further results that
will complete the proof of the Geometrization Conjecture.





As should be clear from the above overview, Perelman's argument did
not arise in a vacuum. Firstly, it resides in a context provided by
the general theory of Riemannian manifolds. In particular,  various
notions of convergence of sequences of manifolds play a crucial
role. The most important is geometric convergence (smooth
convergence on compact subsets). Even more importantly, Perelman's
argument resides in the context of the theory of the Ricci flow
equation, introduced by Richard Hamilton and extensively studied by
him and others. Perelman makes use of almost every previously
established result for $3$-dimensional Ricci flows. One exception is
Hamilton's  proposed classification results for three-dimensional
singularities. These are replaced by Perelman's strong qualitative
description of singularity development for Ricci flows on compact
three-manifolds.

The first five chapters of the book review the necessary background
material from these two subjects. Chapters 6 through 11 then explain
Perelman's advances. In Chapter 12 we introduce the standard
solution, which is the manifold
constructed by hand that one `glues in' in doing surgery. Chapters
13 through 17 describe in great detail the surgery process and prove
the main analytic and topological estimates that are needed to show
that one can continue the process for all positive time. At the end
of Chapter 17 we have established Theorem~\ref{surgery}. Chapter 18
discusses the finite-time extinction result. Chapter 19 is an
appendix on some topological results that were needed in the surgery
analysis in Chapters 13-17.

\section{Background material from Riemannian geometry}

\subsection{Volume and injectivity radius}

One important general concept that is used throughout is the notion
of a manifold being non-collapsed at a point. Suppose that we have a
point $x$ in a complete Riemannian $n$-manifold. Then we say that
the manifold is {\em
$\kappa$-non-collapsed} at $x$
provided that the following holds: For any $r$ such that the norm of
the Riemannian curvature tensor, $|\Rm|$, is $\le r^{-2}$ at
all points of the metric ball, $B(x,r)$, of radius $r$ centered at
$x$,  we have $\Vol\, B(x,r)\ge \kappa r^n$. There is a
relationship between this notion and the injectivity
radius of $M$ at $x$. Namely, if $|\Rm|\le r^{-2}$ on $B(x,r)$ and if $B(x,r)$ is
$\kappa$-non-collapsed then the injectivity radius of $M$ at $x$ is
greater than or equal to  a positive constant that depends only on
$r$ and $\kappa$. The advantage of working with the volume
non-collapsing condition is that, unlike for the injectivity radius,
there is a simple equation for the evolution of volume under Ricci
flow.

Another important general result is the Bishop-Gromov volume
comparison result that says that if the
Ricci curvature of a complete Riemannian $n$-manifold $M$ is bounded
below by a constant $(n-1)K$  then for any $x\in M$ the ratio of the
volume of $B(x,r)$ to the volume of the ball of radius $r$ in the
space of constant curvature $K$ is a non-increasing function whose
limit as $r\rightarrow 0$ is $1$.

All of these basic facts from Riemannian geometry are reviewed in
the first chapter.


\subsection{Manifolds of non-negative curvature}

For reasons that should be clear from the above description and in
any event will become much clearer shortly, manifolds of
non-negative curvature play an extremely important role in the
analysis of Ricci flows with surgery. We need several general
results about them. The first is the soul theorem for manifolds of
non-negative sectional curvature. A {\em soul} is a compact, totally
geodesic submanifold. The entire manifold is diffeomorphic to the
total space of a vector bundle over any of its souls. If a
non-compact $n$-manifold has positive sectional curvature, then any
soul for it is a point, and in particular, the manifold is
diffeomorphic to Euclidean space. In addition, the distance function
$f$ from a soul has the property that for every $t>0$ the pre-image
$f^{-1}(t)$ is homeomorphic to an $(n-1)$-sphere and the pre-image
under this distance function of any non-degenerate interval
$I\subset \Ar ^+$ is homeomorphic to $S^{n-1}\times I$.


Another important result is the splitting theorem, which says that, if a complete manifold of non-negative
sectional curvature has a geodesic line (an isometric copy of $\Ar$)
that is distance minimizing between every pair of its points, then
that manifold is a metric product of a manifold of one lower
dimension and $\Ar$. In particular, if a complete $n$-manifold of
non-negative sectional curvature has two ends then it is a metric
product $N^{n-1}\times \Ar$ where $N^{n-1}$ is a compact manifold.


Also, we need some of the elementary comparison results from Toponogov
theory. These compare ordinary triangles in the
Euclidean plane with triangles in a manifold of non-negative sectional
curvature whose sides are minimizing geodesics in that manifold.



\subsection{Canonical neighborhoods}

Much of the analysis of the geometry of Ricci flows revolves around
the notion of canonical neighborhoods.
Fix some $\epsilon>0$ sufficiently small. There are two types of
non-compact canonical neighborhoods: $\epsilon$-necks
and $\epsilon$-caps. An
$\epsilon$-neck in a Riemannian $3$-manifold $(M,g)$ centered at a
point $x\in M$ is a submanifold $N\subset M$ and a diffeomorphism
$\psi\colon S^2\times (-\epsilon^{-1},\epsilon^{-1})\to N$ such that
$x\in \psi(S^2\times \{0\})$ and such that the pullback of the
rescaled metric, $\psi^*(R(x)g)$, is within $\epsilon$ in the
$C^{[1/\epsilon]}$-topology of the product of the round metric of
scalar curvature $1$ on $S^2$ with the usual metric on the interval
$(-\epsilon^{-1},\epsilon^{-1})$. (Throughout, $R(x)$ denotes the
scalar curvature of $(M,g)$ at the point $x$.) An $\epsilon$-cap is
a non-compact submanifold ${\mathcal C}\subset M$ with the property
that a neighborhood $N$ of infinity in ${\mathcal C}$ is an
$\epsilon$-neck, such that every point of $N$ is the center of an
$\epsilon$-neck in $M$, and such that the {\em
core}, ${\mathcal C}\setminus
\overline N$, of the $\epsilon$-cap is diffeomorphic to either a
$3$-ball or a punctured $\Ar P^3$. It will also be important to
consider $\epsilon$-caps that, after rescaling to make $R(x)=1$ for
some point $x$ in the cap, have bounded geometry (bounded diameter,
bounded ratio of the curvatures at any two points, and bounded
volume). If $C$ represents the bound for these quantities, then we
call the cap an $(C,\epsilon)$-cap. 
 An
$\epsilon$-tube in $M$ is a submanifold of
$M$ diffeomorphic to $S^2\times (0,1)$  which is a union of
$\epsilon$-necks and with the property that each point of the
$\epsilon$-tube is the center of an $\epsilon$-neck in $M$.







There are two other types of canonical neighborhoods in
$3$-manifolds -- (i) a $C$-component and (ii)
an $\epsilon$-round component. The
$C$-component is a compact, connected Riemannian manifold of
positive sectional curvature diffeomorphic to either $S^3$ or $\Ar
P^3$ with the property that rescaling  the metric by $R(x)$ for any
$x$ in the component produces a Riemannian manifold whose  diameter
 is at most $C$, whose sectional curvature at any point and in any
$2$-plane direction is between $C^{-1}$ and $C$, and whose volume is
between $C^{-1}$ and $C$. An $\epsilon$-round component is a
component on which the metric rescaled by $R(x)$ for any $x$ in the
component is within $\epsilon$ in the $C^{[1/\epsilon]}$-topology of
a round metric of scalar curvature one.

As we shall see, the singularities at time $T$ of a $3$-dimensional
Ricci flow are contained in subsets that are unions of canonical
neighborhoods with respect to the metrics at nearby, earlier times
$t'<T$. Thus, we need to understand the topology of manifolds that
are unions of $\epsilon$-tubes and $\epsilon$-caps. The fundamental
observation is that, provided that $\epsilon$ is sufficiently small,
when two $\epsilon$-necks intersect (in more than a small
neighborhood of the boundaries) their product structures almost line
up, so that the two $\epsilon$-necks can be glued together to form a
manifold fibered by $S^2$'s. Using this idea we show that, for
$\epsilon>0$ sufficiently small, if a connected manifold is a union
of $\epsilon$-tubes and $\epsilon$-caps  then it is diffeomorphic to
$\Ar^3$,  $S^2\times \Ar$, $S^3$, $S^2\times S^1$, $\Ar P^3\#\Ar
P^3$,  the total space of a line bundle over $\Ar P^2$, or  the
non-orientable $2$-sphere bundle over $S^1$. This topological result
is proved in the appendix at the end of the book. {\bf We shall fix
$\epsilon>0$ sufficiently small so that these results hold.}

There is one result relating canonical neighborhoods and manifolds
of positive curvature of which we make repeated use: Any complete
$3$-manifold of positive curvature does not admit $\epsilon$-necks
of arbitrarily high curvature. In particular, if $M$ is a complete
Riemannian $3$-manifold with the property that every point of scalar
curvature greater than $r_0^{-2}$ has a canonical neighborhood, then
$M$ has bounded curvature. This turns out to be of central
importance and is used repeatedly.

All of these basic facts about Riemannian manifolds of non-negative
curvature are recalled in the second chapter.

\section{Background material from Ricci flow}

Hamilton \cite{Hamilton3MPRC} introduced the Ricci flow
equation,
$$\frac{\partial g(t)}{\partial t}=-2\Ric(g(t)).$$
This is an evolution equation for a one-parameter family of Riemannian metrics
$g(t)$ on a smooth manifold $M$. The Ricci flow equation is weakly parabolic
and is strictly parabolic modulo the `gauge group', which is the group of
diffeomorphisms of the underlying smooth manifold. One should view this
equation as a non-linear, tensor version of the heat equation. From it, one can
derive the evolution equation for the Riemannian metric tensor, the Ricci
tensor, and the scalar curvature function. These are all parabolic equations.
For example, the evolution equation for scalar curvature $R(x,t)$ is
\begin{equation}\label{scalarflow} \frac{\partial R}{\partial t}(x,t)=\triangle
R(x,t)+2|\Ric(x,t)|^2, \end{equation} illustrating the
similarity with the heat equation. (Here $\triangle$ is the
Laplacian with non-positive spectrum.)

\subsection{First results}

Of course, the first results we need are uniqueness and short-time existence
for solutions to the Ricci flow equation for compact manifolds. These results
were proved by Hamilton (\cite{Hamilton3MPRC}) using the Nash-Moser inverse
function theorem,
 (\cite{HamiltonIFTNM}). These results are standard for
strictly parabolic equations. By now there is a fairly standard method for
working `modulo' the gauge group (the group of diffeomorphisms) and hence
arriving at a strictly parabolic situation where the classical existence,
uniqueness and smoothness results apply. The method for the Ricci flow equation
goes under the name of `DeTurck's trick.'

There is also a result that allows us to patch together local solutions
$(U,g(t)),\ a\le t\le b$, and $(U,h(t)),\ b\le t\le c$, to form a smooth
solution defined on the interval $a\le t\le c$ provided that $g(b)=h(b)$.

Given a Ricci flow $(M,g(t))$ we can always translate, replacing $t$
by $t+t_0$ for some fixed $t_0$, to produce a new Ricci flow. We can
also rescale by any positive constant $Q$ by setting
$h(t)=Qg(Q^{-1}t)$ to produce a new Ricci flow.



\subsection{Gradient shrinking solitons}


Suppose that $(M,g)$ is a complete Riemannian manifold, and suppose
that there is a constant $\lambda>0$ with the property that
$$\Ric(g)=\lambda g.$$
In this case, it is easy to see that there is a Ricci flow given by
$$g(t)=(1-2\lambda t)g.$$ In particular, all the metrics in this flow
differ by a constant factor depending on time and the metric is a
decreasing function of time. These are called {\em shrinking
solitons}. Examples are compact
manifolds of constant positive Ricci curvature.

There is a closely related, but more general, class of examples: the
{\em gradient shrinking solitons}. Suppose that $(M,g)$ is
 a complete Riemannian manifold, and suppose that  there is a constant
$\lambda>0$ and a function $f\colon M\to \Ar$ satisfying
$$\Ric(g)=\lambda g-\Hess^{g}f.$$
 In this case, there is a Ricci flow which is a shrinking
 family after we pull back by the one-parameter family of diffeomorphisms generated by
 the time-dependent vector field $\frac{1}{1-2\lambda t}\nabla_{g} f$.
 An example of a gradient shrinking soliton is the manifold
 $S^2\times \Ar$ with the family of metrics being the product of the
 shrinking family of round metrics on $S^2$ and the constant
 family of standard metrics on $\Ar$. The function $f$ is $s^2/4$ where $s$ is the
 Euclidean parameter on $\Ar$.



\subsection{Controlling higher derivatives of curvature}

Now let us discuss the smoothness results for geometric limits. The
general result along these lines is Shi's theorem, see \cite{Shi1,Shi2}. Again, this is a standard type of
result for parabolic equations. Of course, the situation here is
complicated somewhat by the existence of the gauge group. Roughly,
Shi's theorem says the following. Let us denote by $B(x,t_0,r)$ the
metric ball in $(M,g(t_0))$ centered at $x$ and of radius $r$. If we
can control the norm of the Riemannian curvature tensor on a
backward neighborhood of the form $B(x,t_0,r)\times [0,t_0]$, then
for each $k>0$ we can control the $k^{th}$ covariant derivative of
the curvature on $B(x,t_0,r/2^k)\times [0,t_0]$ by a constant over
$t^{k/2}$. This result has many important consequences in our study
because it tells us that geometric limits are smooth limits. Maybe
the first result to highlight is the fact (established earlier by
Hamilton) that if $(M,g(t))$ is a Ricci flow defined on $0\le
t<T<\infty$, and if the Riemannian curvature is uniformly bounded
for the entire flow, then the Ricci flow extends past time $T$.



In the third chapter this material is reviewed and, where necessary,
slight variants of results and arguments in the literature are
presented.


\subsection{Generalized Ricci flows}

Because we cannot restrict our attention to Ricci flows, but rather
must consider more general objects, Ricci flows with surgery, it is
important to establish the basic analytic results and estimates in a
context more general than that of Ricci flow. We choose to do this
in the context of generalized Ricci flows.

A generalized three-dimensional Ricci flow consists of a smooth
four-dimensional manifold ${\mathcal M}$ (space-time) with a time function ${\bf t}\colon {\mathcal M}\to \Ar$
and a smooth vector field $\chi$. These are required to satisfy:
\begin{enumerate}
\item Each $x\in {\mathcal M}$ has a neighborhood of the form
$U\times J$, where $U$ is an open subset in $\Ar^3$ and $J\subset
\Ar$ is an interval, in which ${\bf t}$ is the projection onto $J$
and $\chi$ is the unit vector field tangent to the one-dimensional
foliation $\{u\}\times J$ pointing in the direction of increasing
${\bf t}$. We call ${\bf t}^{-1}(t)$ the $t$ time-slice. It is a
smooth $3$-manifold.
\item The image ${\bf t}({\mathcal M})$ is a connected interval $I$ in
$\Ar$, possibly infinite. The boundary of ${\mathcal M}$ is the
pre-image under ${\bf t}$ of the boundary of $I$.
\item The level sets ${\bf t}^{-1}(t)$ form a codimension-one foliation of
${\mathcal M}$, called the horizontal foliation, with the boundary
components of ${\mathcal M}$ being leaves.
\item There is a metric $G$ on the horizontal distribution, i.e., the distribution
tangent to the level sets of ${\bf t}$. This metric induces a
Riemannian metric on each $t$ time-slice varying smoothly as we vary
the time-slice. We define the curvature of $G$ at a point
$x\in{\mathcal M}$ to be the curvature of the Riemannian metric
induced by $G$ on the time-slice $M_t$ at $x$.
\item Because of the first property the integral curves of $\chi$
preserve the horizontal foliation and hence the horizontal
distribution. Thus, we can take the Lie derivative of $G$ along
$\chi$. The Ricci flow equation is then
$${\mathcal L}_\chi(G)=-2\Ric(G).$$
\end{enumerate}


Locally in space-time the horizontal metric is simply a smoothly varying family
of Riemannian metrics on a fixed smooth manifold and the evolution equation is
the ordinary Ricci flow equation. This means that the usual formulas for the
evolution of the curvatures as well as much of the analytic analysis of Ricci
flows still hold in this generalized context. In the end, a Ricci flow with
surgery is a more singular type of space-time, but it will have an open dense
subset which is a generalized Ricci flow, and all the analytic estimates take
place in this open subset.

The notion of canonical neighborhoods make sense in the context of
generalized Ricci flows. There is also the notion of a strong
$\epsilon$-neck. Consider an embedding $\psi\colon \left(S^2\times
(-\epsilon^{-1},\epsilon^{-1})\right)\times (-1,0]$ into space-time
such that the time function pulls back to the projection onto
$(-1,0]$ and the vector field $\chi$ pulls back to
$\partial/\partial t$. If there is such an embedding into an
appropriately shifted and rescaled version of the original
generalized Ricci flow so that the pull-back of the rescaled
horizontal metric is within $\epsilon$ in the
$C^{[1/\epsilon]}$-topology of the product of the shrinking family
of round $S^2$'s with the Euclidean metric on
$(-\epsilon^{-1},\epsilon^{-1})$, then we say that $\psi$ is a
strong $\epsilon$-neck in the
generalized Ricci flow.




\subsection{The maximum principle}\label{sectmax}

The Ricci flow equation satisfies various forms of the maximum
principle. The fourth chapter explains this principle, which is due
to Hamilton (see Section 4 of \cite{Hamiltonsurvey}), and derives
many of its consequences,  which are also due to Hamilton (cf.
\cite{HamiltonNSRF3M}). This principle and its consequences are at
the core of all the detailed results about the nature of the flow.
We illustrate the idea by considering the case of the scalar
curvature. A standard scalar maximum principle argument applied to
Equation~(\ref{scalarflow}) proves that the minimum of the scalar
curvature is a non-decreasing function of time. In addition, it
shows that if the minimum of scalar curvature at time $0$ is
positive then we have
$$R_\mathrm{min}(t)\ge R_\mathrm{min}(0)\left(\frac{1}{1-\frac{2t}{n}R_\mathrm{min}(0)}\right),$$ and thus the equation must develop a singularity at or
before time $n/\left(2R_\mathrm{min}(0)\right)$.

While the above result about the scalar curvature is important and
is used repeatedly, the most significant uses of the maximum
principle involve the tensor version, established by Hamilton, which
applies for example to the Ricci tensor and the full curvature
tensor.  These have given the most significant understanding of the
Ricci flows, and they form the core of the arguments that Perelman
uses in his application of Ricci flow to $3$-dimensional topology.
Here are the main results established by Hamilton:

\begin{enumerate}
\item[(1)] For $3$-dimensional flows, if the Ricci curvature is positive,
then the family of metrics becomes singular at finite time and as
the family becomes singular, the metric becomes closer and closer to
round; see \cite{Hamilton3MPRC}.
\item[(2)] For $3$-dimensional flows, as the scalar curvature goes to $+\infty$
the ratio of the absolute value of any negative eigenvalue of the
Riemannian curvature to the largest positive eigenvalue goes to
zero; see \cite{HamiltonNSRF3M}. This condition is called {\em
pinched toward positive curvature}.
\item[(3)] Motivated by a Harnack inequality for the heat equation established
by Li-Yau \cite{LiYau}, Hamilton established a Harnack
inequality for the
curvature tensor under the Ricci flow for complete manifolds
$(M,g(t))$  with bounded, non-negative curvature operator; see
\cite{Hamiltonharnack}. In the applications to three dimensions, we
shall need the following consequence for the scalar curvature:
Suppose that $(M,g(t))$ is  a Ricci flow defined for all $t\in
[T_0,T_1]$ of complete manifolds of non-negative curvature operator
with bounded curvature. Then
$$\frac{\partial R}{\partial t}(x,t)+\frac{R(x,t)}{t-T_0}\ge 0.$$
In particular,
 if $(M,g(t))$ is an ancient solution (i.e.,
defined for all $t\le 0$) of bounded, non-negative curvature then
$\partial R(x,t)/\partial t\ge 0$.
\item[(4)] If a complete $3$-dimensional Ricci flow
$(M,g(t)),\ 0\le t\le T$, has non-negative curvature, if $g(0)$ is
not flat, and if there is at least one point $(x,T)$ such that the
Riemannian curvature tensor of $g(T)$ has a flat direction in
$\wedge^2TM_x$, then $M$ has a cover $\widetilde M$ so that for each
$t>0$ the Riemannian manifold $(\widetilde M,g(t))$ splits as a
Riemannian product of a surface of positive curvature and a
Euclidean line. Furthermore, the flow on the cover $\widetilde M$ is
the product of a $2$-dimensional flow and the trivial
one-dimensional Ricci flow on the line; see Sections 8 and 9 of
\cite{Hamilton4MPCO}.
\item[(5)] In particular, there is no Ricci flow of non-negative curvature tensor
$(U,g(t)),$ defined for $\ 0\le t\le T$ with $T>0$, such that
$(U,g(T))$ is isometric to an open subset in a non-flat,
$3$-dimensional metric cone.
\end{enumerate}


\subsection{Geometric limits}


In the fifth chapter we discuss geometric
limits of Riemannian manifolds and of Ricci flows. Let us review the
history of these ideas. The first results about geometric limits of
Riemannian manifolds go back to Cheeger in his thesis in 1967; see
\cite{Cheeger}. Here Cheeger obtained topological results. In
\cite{Gromov} Gromov proposed that geometric limits should exist in
the Lipschitz topology and suggested a result along these lines,
which also was known to Cheeger.
 In
\cite{GreeneWu}, Greene-Wu gave a rigorous proof of the compactness
theorem suggested by Gromov and also enhanced the convergence to be
$C^{1,\alpha}$-convergence by using harmonic coordinates; see also
\cite{SPeters}. Assuming that all the derivatives of curvature are
bounded, one can apply elliptic theory to the expression of
curvature in harmonic coordinates and deduce $C^\infty$-convergence.
These ideas lead to various types of compactness results  that go
under the name Cheeger-Gromov compactness for Riemannian manifolds.
Hamilton in \cite{Hamiltonlimits} extended these results to Ricci
flows. We shall use the compactness results for both Riemannian
manifolds and for Ricci flows. In a different direction, geometric
limits were extended to the non-smooth context by Gromov in
\cite{Gromov} where he introduced a weaker topology, called the
Gromov-Hausdorff topology and proved a compactness theorem.



 Recall that a
sequence of based Riemannian manifolds $(M_n,g_n,x_n)$ is said to
{\em converge geometrically} to a
based, complete Riemannian manifold $(M_\infty,g_\infty,x_\infty)$
if there is a sequence of open subsets $U_n\subset M_\infty$ with
compact closures, with $x_\infty\in U_1\subset \overline U_1\subset
U_2\subset \overline U_2\subset U_3\subset \cdots $ with
$\cup_nU_n=M_\infty$, and embeddings $\varphi_n\colon U_n\to M_n$
sending $x_\infty$ to $x_n$ so that the pull back metrics,
$\varphi_n^*g_n$, converge uniformly on compact subsets of
$M_\infty$ in the $C^\infty$-topology to $g_\infty$. Notice that the
topological type of the limit can be different from the topological
type of the manifolds in the sequence. There is a similar notion of
geometric convergence for a sequence of based Ricci flows.

Certainly, one of the most important consequences of Shi's results,
cited above, is that, in concert with  Cheeger-Gromov  compactness,
it allows us to form smooth geometric limits of sequences of based
Ricci flows. We have the following result of Hamilton's; see
\cite{Hamiltonlimits}:

\begin{theorem}[Compactness theorem for based Ricci flows (Hamilton)]
\label{thm:compactness-based-ricci-flows}
Suppose we have a sequence of based Ricci flows
$(M_n,g_n(t),(x_n,0))$ defined for $t\in (-T,0]$ with the
$(M_n,g_n(t))$ being complete. Suppose that:
\begin{enumerate}
\item[(1)] There is $r>0$ and $\kappa>0$ such that for every $n$ the metric ball $B(x_n,0,r)\subset
(M_n,g_n(0))$ is $\kappa$-non-collapsed.
\item[(2)] For each $A<\infty$ there is $C=C(A)<\infty$ such that the Riemannian
curvature on $B(x_n,0,A)\times (-T,0]$ is bounded by $C$.
\end{enumerate}
Then after passing to a subsequence there is a geometric limit which
is a based Ricci flow $(M_\infty,g_\infty(t),(x_\infty,0))$ defined
for $t\in (-T,0]$.
\end{theorem}

To emphasize, the two conditions that we must check in order to
extract a geometric limit of a subsequence based at points at time
zero are: (i) uniform non-collapsing at the base point in the time
zero metric, and (ii) for each $A<\infty$ uniformly bounded
curvature for the restriction of the flow to the metric balls of
radius $A$ centered at the base points.

Most steps in Perelman's argument require invoking this result in order to form
limits of appropriate sequences of Ricci flows, often rescaled to make the
scalar curvatures at the base point equal to $1$. If, before rescaling, the
scalar curvature at the base points goes to infinity as we move through the
sequence, then the resulting limit of the rescaled flows has non-negative
sectional curvature. This is a consequence of the fact that the sectional
curvatures of the manifolds in the sequence are uniformly pinched toward
positive. It is for exactly this reason that non-negative curvature plays such
an important role in the study of singularity development in three-dimensional
Ricci flows.




\section{Perelman's advances}

So far we have been discussing the results that were known before
Perelman's work. They concern almost exclusively Ricci flow (though
Hamilton in \cite{Hamilton4MPIC} had introduced the notion of
surgery and proved that surgery can be performed preserving the
condition that the curvature is pinched toward positive, as in (2)
above). Perelman extended in two essential ways the analysis of
Ricci flow -- one involves the introduction of a new analytic
functional, {\em the reduced length}, which is the tool by which he
establishes the needed non-collapsing results, and the other is a
delicate combination of geometric limit ideas and consequences of
the maximum principle together with the non-collapsing results in
order to establish bounded curvature at bounded distance results.
These are used to prove in an inductive way the existence of
canonical neighborhoods, which is a crucial ingredient in proving
that is possible to do surgery iteratively, creating a flow defined
for all positive time.

While it is easiest to formulate and consider these techniques in
the case of Ricci flow, in the end one needs them in the more
general context of Ricci flow with surgery since we inductively
repeat the surgery process, and in order to know at each step that
we can perform surgery we need to apply these results to the
previously constructed Ricci flow with surgery. We have chosen to
present these new ideas only once -- in the context of generalized
Ricci flows -- so that we can derive the needed consequences in all
the relevant contexts from this one source.

\subsection{The reduced length function}

In Chapter~\ref{lengthfn} we come to
the first of Perelman's major contributions. Let us first describe
it in the context of an ordinary three-dimensional Ricci flow, but
viewing the Ricci flow as a horizontal metric on a space-time which
is the manifold $M\times I$, where $I$ is the interval of definition
of the flow. Suppose that $I=[0,T)$ and fix $(x,t)\in M\times
(0,T)$. We consider paths $\gamma(\tau),\ 0\le \tau\le
\overline\tau$, in space-time with the property that for every
$\tau\le \overline\tau$ we have $\gamma(\tau)\in M\times \{t-\tau\}$
and $\gamma(0)=x$. These paths are said to be {\em parameterized by
backward time}.  The ${\mathcal
L}$-{\em length} of such a path is given by
$${\mathcal L}(\gamma)=\int_0^{\overline
\tau}\sqrt{\tau}\left(R(\gamma(\tau))+|\gamma'(\tau)|^2\right)d\tau,$$
where the derivative on $\gamma$ refers to the spatial derivative.
There is also the closely related {\em reduced length}
$$\ell(\gamma)=\frac{{\mathcal L}(\gamma)}{2\sqrt{\overline\tau}}.$$
There is a theory for the functional ${\mathcal L}$ analogous to the
theory for the usual energy function\footnote{Even though this
functional is called a length, the presence of the
$|\gamma'(\tau)|^2$ in the integrand means that it behaves more like
the usual energy functional for paths in a Riemannian manifold.}.
 In
particular, there is the notion of an ${\mathcal
L}$-geodesic, and the reduced length
as a function on space-time $\ell_{(x,t)}\colon M\times [0,t)\to
\Ar$. One establishes a crucial monotonicity for this reduced length along ${\mathcal
L}$-geodesics. Then one defines the reduced volume
$$\widetilde
V_{(x,t)}(U\times\{\bar t\})=\int_{U\times\{\bar
t\}}\bar\tau^{-3/2}e^{-\ell_{(x,t)}(q,\bar\tau)}d\vol_{g(\bar\tau}(q),$$ where, as before $\bar\tau=t-\bar t$.
Because of the monotonicity of $\ell_{(x,t)}$ along ${\mathcal
L}$-geodesics, the reduced volume is also
non-increasing under the flow (forward in $\bar\tau$ and hence
backward in time) of open subsets along ${\mathcal L}$-geodesics.
This is the fundamental tool which is used to establish
non-collapsing results which in turn are essential in proving the
existence of geometric limits.







The definitions and the analysis of the reduced length function and
the reduced volume as well as the monotonicity results are valid in
the context of the generalized Ricci flow. The only twist to be
aware of is that in the more general context one cannot always
extend ${\mathcal L}$-geodesics; they may run `off the edge' of
space-time. Thus, the reduced length function and reduced volume
cannot be defined globally, but only on appropriate open subsets of
a time-slice (those reachable by minimizing ${\mathcal
L}$-geodesics). But as long as one can flow an open set $U$ of a
time-slice along minimizing ${\mathcal L}$-geodesics in the
direction of decreasing $\bar\tau$, the reduced volumes of the
resulting family of open sets form a monotone non-increasing
function of $\bar\tau$. This turns out to be sufficient to extend
the non-collapsing results to Ricci flow with surgery provided that
we are careful in how we choose the parameters that go into the
definition of the surgery process.

\subsection{Application to non-collapsing results}

As we indicated in the previous paragraph, one of the main
applications of the reduced length function is to prove
non-collapsing results for three-dimensional Ricci flows with
surgery. In order to make this  argument work, one takes a weaker
notion of $\kappa$-non-collapsed by
making a stronger curvature bound assumption: one considers
 points $(x,t)$ and constants $r$ with the property that $|\Rm|\le r^{-2}$ on
$P(x,t,r,-r^2)=B(x,t,r)\times (t-r^2,t]$. The
$\kappa$-non-collapsing condition applies to these balls and says
that $\Vol(B(x,t,r))\ge \kappa r^3.$ The basic idea in proving
non-collapsing is to use the fact that as we flow forward in time
via minimizing ${\mathcal L}$-geodesics the reduced volume is a
non-decreasing function. Hence, a lower bound of the reduced volume
of an open set at an earlier time implies the same lower bound for
the corresponding open subset at a later time. This is contrasted
with direct computations (related to the heat kernel in $\Ar^3$)
that say if the manifold is highly collapsed near $(x,t)$ (i.e.,
satisfies the curvature bound above but is not
$\kappa$-non-collapsed for some small $\kappa$) then the reduced
volume $\widetilde V_{(x,t)}$ is small at times close to $t$. Thus,
to show that the manifold is non-collapsed at $(x,t)$ we need only
find an open subset at an earlier time that is reachable by
minimizing ${\mathcal L}$-geodesics and that has a reduced volume
bounded away from zero.




One case where it is easy to do this is when we have a Ricci flow of
compact manifolds or of complete manifolds of non-negative
curvature. Hence, these manifolds are non-collapsed at all points
with a  non-collapsing constant that depends only on the geometry of
the initial metric of the Ricci flow. Non-collapsing results are
crucial and are used repeatedly in dealing with Ricci flows with
surgery in Chapters 10 -- 17, for these give one of the two
conditions required in order to take geometric limits.


\subsection{Application to ancient $\kappa$-non-collapsed solutions}

There is another important application of the length function, which
is to the study of non-collapsed, ancient solutions in dimension three. In the case that the generalized Ricci
flow is an ordinary Ricci flow either on a compact manifold or on a
complete manifold (with bounded curvatures) one can say much more
about the reduced length function and the reduced volume. Fix a
point $(x_0,t_0)$ in space-time. First of all, one shows that every
point $(x,t)$ with $t<t_0$ is reachable by a minimizing ${\mathcal
L}$-geodesic and thus that the reduced length is defined as a
function on all points of space at all times $t<t_0$. It turns out
to be a locally Lipschitz function in both space and time and hence
its gradient and its time derivative exist as $L^2$-functions and
satisfy important differential inequalities in the weak sense.

These results apply to a class of Ricci flows called
$\kappa$-solutions, where $\kappa$ is a
positive constant. By definition a $\kappa$-solution is a Ricci flow
defined for all $t\in (-\infty,0]$, each time-slice is a non-flat,
complete $3$-manifold of non-negative, bounded curvature and each
time-slice is $\kappa$-non-collapsed. The differential inequalities
for the reduced length from any point $(x,0)$ imply that, for any
$t<0$, the minimum value of $\ell_{(x,0)}(y,t)$ for all $y\in M$ is
at most $3/2$. Furthermore, again using the differential
inequalities for the reduced length function, one shows that for any
sequence $t_n\rightarrow -\infty$, and any points $(y_n,t_n)$ at
which the reduced length function is bounded above by $3/2$, there
is a subsequence of based Riemannian manifolds,
$(M,\frac{1}{|t_n|}g(t_n),y_n)$, with a geometric limit, and this
limit is a gradient shrinking soliton. This gradient shrinking
soliton is called an {\em asymptotic
soliton} for the
original $\kappa$-solution, 




The point is that there are only two types of gradient shrinking
solitons in dimension three -- (i) those finitely covered by a
family of shrinking round $3$-spheres and (ii) those finitely
covered by a family of shrinking round cylinders $S^2\times \Ar$. If
a $\kappa$-solution has a gradient shrinking soliton of the first
type then it is in fact isomorphic to its gradient shrinking
soliton. More interesting is the case when the $\kappa$-solution has
a gradient shrinking soliton which is of the second type. If the
$\kappa$-solution does not have strictly positive curvature, then it
is isomorphic to its gradient shrinking soliton. Furthermore, there
is a constant $C_1<\infty$ depending on $\epsilon$ (which remember
is taken sufficiently small) such that a $\kappa$-solution of
strictly positive curvature is either a $C_1$-component, or is a
union of cores of $(C_1,\epsilon)$-caps and points that are the
center points of $\epsilon$-necks.

In order to prove the above results (for example the uniformity of $C_1$ as
above over all $\kappa$-solutions)  one needs the following result:
\begin{theorem}[Compactness of the space of based kappa-solutions]
\label{thm:compactness-kappa-solutions}
The space of based $\kappa$-solutions,  based at points $(x,0)$ with
$R(x,0)=1$, is compact.
\end{theorem}

This result does not generalize to ancient solutions that are not
non-collapsed because, in order to prove compactness, one has to
take limits of subsequences, and in doing this the non-collapsing
hypothesis is essential. See Hamilton's work \cite{Hamiltonsurvey}
for more on general ancient solutions (i.e., those that are not
necessarily non-collapsed).

Since $\epsilon>0$ is sufficiently small so that all the results from the
appendix about manifolds covered by $\epsilon$-necks and $\epsilon$-caps hold,
the above results about gradient shrinking solitons lead to a rough qualitative
description of all $\kappa$-solutions. There are those which do not have
strictly positive curvature. These are gradient shrinking solitons, either an
evolving family of round $2$-spheres times $\Ar$ or the quotient of this family
by an involution. Non-compact $\kappa$-solutions of strictly positive curvature
are diffeomorphic to $\Ar^3$ and are the union of an $\epsilon$-tube and a core
of a $(C_1,\epsilon)$-cap. The compact ones of strictly positive curvature are
of two types. The first type are positive, constant curvature shrinking
solitons. Solutions of the second type are diffeomorphic to either $S^3$ or
$\Ar P^3$. Each time-slice of a $\kappa$-solution of the second type is either
of uniformly bounded geometry (curvature, diameter, and volume) when rescaled
so that the scalar curvature at a point is one, or admits an $\epsilon$-tube
whose complement is either a disjoint union of the cores of two
$(C_1,\epsilon)$-caps.


This gives a rough qualitative understanding of  $\kappa$-solutions. Either
they are round, or they are finitely covered by the product of a round surface
and a line, or they are a union of $\epsilon$-tubes and cores of
$(C_1,\epsilon)$-caps , or they are diffeomorphic to $S^3$ or $\Ar P^3$ and
have bounded geometry (again after rescaling so that there is a point of scalar
curvature $1$). This is the source of canonical neighborhoods for Ricci flows:
 the point is that this qualitative result remains true for any point $x$ in a
Ricci flow that has an appropriate size neighborhood within
$\epsilon$ in the $C^{[1/\epsilon]}$-topology of a neighborhood in a
$\kappa$-solution. For example, if we have a sequence of based
generalized flows $({\mathcal M}_n,G_n,x_n)$ converging to a based
$\kappa$-solution, then for all $n$ sufficiently large $x$ will have
a canonical neighborhood, one that is either an $\epsilon$-neck
centered at that point, a $(C_1,\epsilon)$-cap whose core contains
the point, a $C_1$-component, or an $\epsilon$-round component.

\subsection{Bounded curvature at bounded distance}

Perelman's other major breakthrough is his result establishing
bounded curvature at bounded distance for blow-up limits of
generalized Ricci flows. As we have alluded to several times, many
steps in the argument require taking (smooth) geometric limits of a
sequence of based generalized flows about points of curvature
tending to infinity.  To study  such a sequence we rescale each term
in the sequence so that its curvature at the base point becomes one.
Nevertheless, in taking such limits we face the problem that even
though the curvature at the point we are focusing on (the points we
take as base points) was originally large and has been  rescaled to
be one, there may be other points in the same time-slice of much
larger curvature, which, even after the rescalings, can tend to
infinity.
 If
these points are  at uniformly bounded (rescaled) distance from the
base points, then they would preclude the existence of a smooth
geometric limit of the based, rescaled flows. In his arguments,
Hamilton avoided this problem by always focusing on points of
maximal curvature (or almost maximal curvature). That method will
not work in this case. The way to deal with this possible problem is
to show that a generalized Ricci flow satisfying appropriate
conditions satisfies the following. For each $A<\infty$ there are
constants $Q_0=Q_0(A)<\infty$ and $Q(A)<\infty$ such that any point
$x$ in such a generalized flow for which the scalar curvature
$R(x)\ge Q_0$ and for any $y$ in the same time-slice as $x$ with
$d(x,y)<AR(x)^{-1/2}$ satisfies $R(y)/R(x)<Q(A)$. As we shall see,
this and the non-collapsing result are the fundamental tools that
allow Perelman to study neighborhoods of points of sufficiently
large curvature by taking smooth limits of rescaled flows, so
essential in studying the prolongation of Ricci flows with surgery.

The basic idea in proving this result is to assume the contrary and
take an incomplete geometric limit of the rescaled flows based at
the counterexample points. The existence of points at bounded
distance with unbounded, rescaled curvature means that there is a
point at infinity at finite distance from the base point where the
curvature blows up. A neighborhood of this point at infinity is
cone-like in a manifold of non-negative curvature. This contradicts
Hamilton's maximum principle result (5) in Chapter~\ref{sectmax})
that the result of a Ricci flow of manifolds of non-negative
curvature is never an open subset of a cone. (We know that any
`blow-up limit' like this has non-negative curvature because of the
curvature pinching result.) This contradiction establishes the
result.



\section{The standard solution and the surgery process}

Now we are ready to discuss three-dimensional Ricci flows with surgery.

\subsection{The standard solution}

In preparing the way for defining the surgery process, we must
construct a metric on the $3$-ball that we shall glue in when we
perform surgery. This we do in Chapter~\ref{stdsolnsect}. We fix a
non-negatively curved, rotationally symmetric metric on $\Ar^3$ that
is isometric near infinity to $S^2\times [0,\infty)$ where the
metric on $S^2$ is the round metric of scalar curvature $1$, and
outside this region has positive sectional curvature,  Any such metric will suffice for the gluing
process, and we fix one and call it the {\em standard
metric}. It is important to understand Ricci flow
with the standard metric as initial metric.  Because of the special
nature of this metric (the rotational symmetry and the asymptotic
nature at infinity), it is fairly elementary to show that there is a
unique solution of bounded curvature on each time-slice to the Ricci
flow equation with the standard metric as the initial metric; this
flow is defined for $0\le t<1$; and for any $T<1$ outside of a
compact subset $X(T)$ the restriction of the flow to $[0,T]$ is
close to the evolving round cylinder. Using the length function, one
shows that the Ricci flow is non-collapsed, and that the bounded
curvature and bounded distance result applies to it.
 This allows one to prove that every point $(x,t)$ in this flow has one of the
following types of neighborhoods:
\begin{enumerate}
\item $(x,t)$ is contained in the core of a $(C_2,\epsilon)$-cap,
where $C_2<\infty$ is a given universal constant depending only on $\epsilon$.
\item $(x,t)$ is the center of a strong $\epsilon$-neck.
\item $(x,t)$ is the center of an evolving $\epsilon$-neck whose
initial slice is at time zero.
\end{enumerate}

These form the second source of models for canonical neighborhoods in a Ricci
flow with surgery. Thus, we shall set $C=C(\epsilon)=\mathrm{max}(C_1(\epsilon),C_2(\epsilon))$ and we shall find $(C,\epsilon)$-canonical
neighborhoods in Ricci flows with surgery.







\subsection{Ricci flows with surgery}

Now it is time to introduce the notion of a Ricci flow with
surgery. To do this we formulate an
appropriate notion of $4$-dimensional space-time that allows for the
surgery operations. We define {\em space-time} to be a
$4$-dimensional Hausdorff singular space with a time function ${\bf
t}$ with the property that each time-slice is a compact, smooth
$3$-manifold, but level sets at different times are not necessarily
diffeomorphic. Generically space-time is a smooth $4$-manifold, but
there are {\em exposed regions} at a discrete
set of times. Near a point in the exposed region  space-time is a
$4$-manifold with boundary. The singular points of space-time are
the boundaries of the exposed regions. Near these, space-time is
modeled on the product of $\Ar^2$ with the square $(-1,1)\times
(-1,1)$, the latter having a topology in which the open sets are, in
addition to the usual open sets, open subsets of $(0,1)\times
[0,1)$,  There is a natural
notion of smooth functions on space-time. These are smooth in the
usual sense on the open subset of non-singular points. Near the
singular points, and in the local coordinates described above, they
are required to be pull-backs from smooth functions on $\Ar^2\times
(-1,1)\times (-1,1)$ under the natural map. Space-time is equipped with a smooth vector field $\chi$
with $\chi({\bf t})=1$.




A Ricci flow with surgery is a smooth horizontal metric $G$ on a
space-time  with the property that the restriction of $G$, ${\bf t}$
and $\chi$ to the open subset of smooth points forms a generalized
Ricci flow.  We call this the {\em associated generalized Ricci
flow} for the Ricci flow with surgery.


\subsection{The inductive conditions necessary for doing surgery}

With all this preliminary work out of the way, we are ready to show
that one can construct Ricci flow with surgery which is precisely
controlled both topologically and metrically. This result is proved
inductively, one interval of time after another, and it is important
to keep track of various properties as we go along to ensure that we
can continue to do surgery. Here we discuss the conditions we verify
at each step.

Fix $\epsilon>0$ sufficiently small and let $C=\mathrm{max}(C_1,C_2)<\infty$, where $C_1$ is the constant associated to
$\epsilon$ for $\kappa$-solutions and $C_2$ is the constant
associated to $\epsilon$ for the standard solution. We say that a
point $x$ in a generalized Ricci flow has a $(C,\epsilon)$-canonical
neighborhood if one of the following holds:
\begin{enumerate}
\item $x$ is contained in a connected component of a
time-slice that is a $C$-component.
\item $x$ is contained in a connected component of its time-slice
that is within $\epsilon$ of round in the
$C^{[1/\epsilon]}$-topology.
\item $x$ is contained in the core of a $(C,\epsilon)$-cap.
\item $x$ is the center of a strong $\epsilon$-neck.
\end{enumerate}


We shall study  Ricci flows with surgery defined for $0\le
t<T<\infty$ whose associated generalized Ricci flows satisfy the
following properties:
\begin{enumerate}
\item The initial metric is normalized, meaning that for the
metric at time zero the norm the Riemann curvature is bounded above
by one and the volume of any ball of radius one is at least half the
volume of the unit ball in Euclidean space.
\item The curvature of the flow is pinched toward positive.
\item There is $\kappa>0$ so that the associated generalized Ricci flow is
$\kappa$-non-collapsed on scales at most $\epsilon$, in the sense
that we require only that balls of radius $r\le \epsilon$ be
$\kappa$-non-collapsed.
\item There is $r_0>0$ such that any point of space-time at which
the scalar curvature is $\ge r_0^{-2}$ has an
$(C,\epsilon)$-canonical neighborhood.
\end{enumerate}

The main result is that, having a Ricci flow with surgery defined on
some time interval satisfying these conditions, it is possible to
extend it to a longer time interval in such a way that it still
satisfies the same conditions, possibly allowing the constants
$\kappa$ and $r_0$ defining these conditions to get closer to zero,
but keeping them bounded away from $0$ on each compact time
interval. We repeat this construction inductively, and, since it is
easy to see that on any compact time interval there can only be a
bounded number of surgeries. In the end, we create a Ricci flow with
surgery defined for all positive time. As far as we know, it may be
the case that in the entire flow defined all the way to infinity
there are infinitely many surgeries.




\subsection{Surgery}

Let us describe how we extend a Ricci flow with surgery satisfying
all the conditions listed above and becoming singular at time
$T<\infty$. Fix $T^-<T$ so that there are no surgery  times in the
interval $[T^-,T)$. Then we can use the Ricci flow to identify all
the time-slices $M_t$ for $t\in [T^-,T)$, and hence view this part
of the Ricci flow with surgery as an ordinary Ricci flow. Because of
the canonical neighborhood assumption, there is an open subset
$\Omega\subset M_{T^-}$ on which the curvature stays bounded as
$t\rightarrow T$. Hence, by Shi's results, there is a limiting
metric at time $T$ on $\Omega$. Furthermore, the scalar curvature is
a proper function, bounded below, from $\Omega$ to $\Ar$, and each
end of $\Omega$ is an $\epsilon$-tube where the cross-sectional area
of the $2$-spheres goes to zero as we go to the end of tube. We call
such tubes {\em $\epsilon$-horns}. We are interested in
$\epsilon$-horns whose boundary is contained in the part of $\Omega$
where the scalar curvature is bounded above by some fixed finite
constant $\rho^{-2}$. We call this region $\Omega_\rho$. Using the
bounded curvature at bounded distance result and using the
non-collapsing hypothesis, one shows that given any $\delta>0$ there
is $h=h(\delta,\rho,r_0)$ such that for any $\epsilon$-horn
${\mathcal H}$ whose boundary lies in $\Omega_\rho$ and for any
$x\in{\mathcal H}$ with $R(x)\ge h^{-2}$, the point $x$ is the
center of a strong $\delta$-neck.

Now we are ready to describe the surgery procedure. It depends on
our choice of standard solution on $\Ar ^3$ and on a choice of
$\delta>0$ sufficiently small. For each $\epsilon$-horn in $\Omega$
whose boundary is contained in $\Omega_\rho$ fix a point of
curvature $(h(\delta, \rho,r_0))^{-2}$ and fix a strong
$\delta$-neck centered at this point. Then we cut the
$\epsilon$-horn open along the central $2$-sphere $S$ of this neck
and remove the end of the $\epsilon$-horn that is cut off by $S$.
Then we glue in a ball of a fixed radius around the tip from the
standard solution, after scaling the metric on this ball by
$(h(\delta,\rho,r_0))^2$. To glue these two metrics together we must
use a partition of unity near the $2$-spheres that are matched.
There is also a delicate point that we first bend in the metrics
slightly so as to achieve positive curvature near where we are
gluing. This is an idea due to Hamilton, and it is needed in order
to show that the condition of curvature pinching toward positive  is
preserved. In addition, we remove all components of $\Omega$ that do
not contain any points of $\Omega_\rho$.


This operation produces a new compact $3$-manifold.  One continues
the Ricci flow with surgery by letting this Riemannian manifold at
time $T$ evolve under the Ricci flow. 




  

\subsection{Topological effect of surgery}

Looking at the
situation just before the surgery time, we see a finite number of
disjoint submanifolds, each diffeomorphic to either $S^2\times I$ or
the $3$-ball, where the curvature is large. In addition there may be
entire components of where the scalar curvature is large. The effect
of $2$-sphere surgery is to do a finite number of ordinary
topological surgeries along $2$-spheres in the $S^2\times I$. This
simply effects a partial connected-sum decomposition and may
introduce new components diffeomorphic to $S^3$. We also remove
entire components, but these are covered by $\epsilon$-necks and
$\epsilon$-caps  so that they have standard topology (each one is
diffeomorphic to $S^3$, $\Ar P^3$, $\Ar P^3\# \Ar P^3$, $S^2\times
S^1$, or the non-orientable $2$-sphere bundle over $S^1$). Also, we
remove $C$-components and $\epsilon$-round components (each of these
is either diffeomorphic to $S^3$ or $\Ar P^3$ or admits a metric of
constant positive curvature). Thus, the topological effect of
surgery is to do a finite number of ordinary $2$-sphere topological
surgeries and to remove a finite number of topologically standard
components.



\section{Extending Ricci flows with surgery}

Now we come to the crux of the argument. We must show that if we
have a Ricci flow with surgery defined for some time $0\le
t<T<\infty$ satisfying the four conditions: normalized initial
metric, curvature pinched toward positive, all points of scalar
curvature $\ge r^{-2}$ have canonical neighborhoods, and the flow is
$\kappa$-non-collapsed on scales $\le \epsilon$; then it is possible
to extend to a Ricci flow with surgery defined past $T$ to a larger
time keeping all these conditions satisfied (possibly with different
constants $r_0'<r_0$ and $\kappa'<\kappa$). In order to do this we
need to choose the surgery parameter $\delta>0$ sufficiently small.
There is also the issue of whether the surgery times can accumulate.

Of course, the initial metric does not change as we extend surgery
so that the condition that the normalized initial metric is clearly
preserved as we extend surgery. As we have already remarked,
Hamilton had proved earlier that one can do surgery in such a way as
to preserve the condition that the curvature is pinched toward
positive. The other two conditions require more work, and, as we
indicated above, the constants may decay to zero as we extend the
Ricci flow with surgery.

If we have all the conditions for the Ricci flow with surgery up to time $T$,
then the analysis of the open subset on which the curvature remains bounded
holds, and given $\delta>0$ sufficiently small, we do surgery on the central
$S^2$ of a strong $\delta$-neck in each $\epsilon$-horn meeting $\Omega_\rho$.
In addition we remove entirely all components that do not contain points of
$\Omega_\rho$. We then glue in the cap from the standard solution. This gives
us a new compact $3$-manifold and we restart the flow from this manifold.

The $\kappa$-non-collapsed result is extended to the new part of the
Ricci flow with surgery using the fact that it holds at times
previous to $T$. To establish this extension one uses ${\mathcal
L}$-geodesics in the associated generalized Ricci flow and reduced
volume as indicated before. In order to get this argument to work,
one must require $\delta>0$ to be sufficiently small; how small is
determined by $r_0$.

The other thing that we must establish is the existence of canonical
neighborhoods for all points of sufficiently large scalar curvature.
Here the argument is by contradiction. We consider all Ricci flows
with surgery that satisfy all four conditions on $[0,T)$ and we
suppose that we can find a sequence of such containing points
(automatically at times $T'>T$) of arbitrarily large curvature where
there are not canonical neighborhoods. In fact, we take the points
at the first such time violating this condition. We base our flows
at these points. Now we consider rescaled versions of the
generalized flows so that the curvature at these base points is
rescaled to one. We are in a position to apply the bounded curvature
and bounded distance results to this sequence, and of course the
$\kappa$-non-collapsing results which have already been established.
There are two possibilities. The first is that the rescaled sequence
converges to an ancient solution. This ancient solution has
non-negative curvature by the pinching hypothesis. General results
about three-manifolds  of non-negative curvature imply that it also
has bounded curvature. It is $\kappa$-non-collapsed. Thus, in this
case the limit is a $\kappa$-solution. This produces the required
canonical neighborhoods for the base points of the tail of the
sequence modeled on the canonical neighborhoods of points in a
$\kappa$-solution. This contradicts the assumption that none of
these points has a canonical neighborhood.


The other possibility is that one can take a partial smooth limit
but that this limit does not extend all the way back to $-\infty$.
The only way this can happen is if there are surgery caps that
prevent extending the limit back to $-\infty$. This means that the
base points in our sequence are all within a fixed distance and time
(after the rescaling) of a surgery region. But in this case results
from the nature of the standard solution show that if we have taken
$\delta>0$ sufficiently small, then the base points have canonical
neighborhoods modeled on the canonical neighborhoods in the standard
solution, again contradicting our assumption that none of the base
points has a canonical neighborhood. In order to show that our base
points have neighborhoods near those of the standard solution, one
appeals to a geometric limit argument as $\delta\rightarrow 0$. This
argument uses the uniqueness of the Ricci flow for the standard
solution. (Actually, Bruce Kleiner pointed out to us that one only
needs a compactness result for the space of all Ricci flows with the
standard metric as initial metric, not uniqueness, and the
compactness result can be proved by the same arguments that prove
the compactness of the space of $\kappa$-solutions.)



Interestingly enough, in order to establish the uniqueness of the
Ricci flow for the standard solution, as well as to prove that this
flow is defined for time $[0,1)$ and to prove that at infinity it is
asymptotic to an evolving cylinder requires the same results --
non-collapsing and the bounded curvature at bounded distance that we
invoked above. For this reason, we order the material described here
as follows. First, we introduce  generalized Ricci flows, and then
introduce the length function in this context and establish the
basic monotonicity results. Then we have a chapter on stronger
results for the length function in the case of complete manifolds
with bounded curvature. At this point we are in a position to prove
the needed results about the Ricci flow from the standard solution.
Then we are ready to define the surgery process and prove the
inductive non-collapsing results and the existence of canonical
neighborhoods.



In this way, one establishes the existence of canonical neighborhoods. Hence,
one can continue to do surgery, producing a Ricci flow with surgery defined for
all positive time.  Since these arguments are inductive, it turns out that the
constants in the non-collapsing and in the canonical neighborhood statements
decay in a predetermined rate as time goes to infinity.


Lastly, there is the issue of ruling out the possibility that the surgery times
accumulate. The idea here is very simple: Under Ricci flow during an elapsed
time $T$, volume increases at most by a  multiplicative factor which is a fixed
exponential of the time $T$. Under each surgery there is a removal of at least
a fixed positive amount of volume depending on the surgery scale $h$, which in
turns depends on $\delta$ and $r_0$. Since both $\delta$ and $r_0$ are bounded
away from zero on each finite interval, there can be at most finitely many
surgeries in each finite interval. Notice that this argument allows for the
possibility that in the entire flow all the way to infinity there are
infinitely many surgeries. It is still unknown whether that possibility ever
happens.

This completes our outline of the proof of Theorem~\ref{surgery}.



\section{Finite-time extinction}

The last topic we discuss is the  proof of the finite-time
extinction for Ricci flows with
initial metrics satisfying the hypothesis  of
Theorem~\ref{finiteext}.


As we present it, the finite extinction result has two steps. The
first step is to show that there is $T<\infty$ (depending on the
initial metric) such that for all $t\ge T$, all connected components
of the $t$-time-slice $M_t$ have trivial $\pi_2$.  First, an easy
topological argument shows that only finitely many of the $2$-sphere
surgeries in a Ricci flow with surgery can be along homotopically
non-trivial $2$-spheres. Thus, after some time $T_0$ all $2$-sphere
surgeries are along homotopically trivial $2$-spheres. Such a
surgery does not affect $\pi_2$. Thus, after time $T_0$, the only
way that $\pi_2$ can change is by removal of components with
non-trivial $\pi_2$. (An examination of the topological types of
components that are removed shows that there are only two types of
such components with non-trivial $\pi_2$: $2$-sphere bundles over
$S^1$ and $\Ar P^3\#\Ar P^3$.) We suppose that at every $t\ge T_0$
there is a component of $M_t$ with non-trivial $\pi_2$. Then we can
find a connected open subset ${\mathcal X}$ of ${\bf
t}^{-1}([T_0,\infty))$ with the property that for each $t\ge  T_0$
the intersection ${\mathcal X}(t)= {\mathcal X}\cap M_t$ is a
component of $M_t$ with non-trivial $\pi_2$. We define a function
$W_2\colon [T_0,\infty)\to \Ar$ associated with such an ${\mathcal
X}$. The value $W_2(t)$ is the minimal area of all homotopically
non-trivial $2$-spheres mapping into ${\mathcal X}(t)$. This minimal
area $W_2(t)$ is realized by a harmonic map of $S^2$ into ${\mathcal
X}(t)$. The function $W_2$ varies continuously under Ricci flow and
at a surgery is lower semi-continuous. Furthermore, using an idea
that goes back to Hamilton (who applied it to minimal disks) one
shows that the forward difference quotient of the minimal area
satisfies
$$\frac{dW_2(t)}{dt}\le -4\pi +\frac{3}{(4t+1)}W_2(t).$$
(Here, the explicit form of the bound for the forward difference
quotient depends on the way we have chosen to normalize initial
metric and also on Hamilton's curvature pinching result.)

But any function $W_2(t)$ with these properties and defined for all
$t>T_0$,  becomes negative at some finite $T_1$ (depending on the
initial value). This is absurd since $W_2(t)$ is the minimum of
positive quantities. This contradiction shows that such a path of
components with non-trivial $\pi_2$ cannot exist for all $t\ge T_0$.
In fact, it even gives a computable upper bound on how long such a
component ${\mathcal X}$, with every time-slice having non-trivial
$\pi_2$, can exist in terms of the minimal area of a homotopically
non-trivial $2$-sphere mapping into ${\mathcal X}(T_0)$. It follows
that there is $T<\infty$ with the property that every component of
$M_T$ has trivial $\pi_2$. This condition then persists for all
$t\ge T$.

Three remarks are in order. This argument showing that eventually
every component of the time-slice $t$ has trivial $\pi_2$ is not
necessary for the topological application (Theorem~\ref{finiteext}),
or indeed, for any other topological application. The reason is the
sphere theorem (see \cite{Hempel}), which says that if $\pi_2(M)$ is
non-trivial then either $M$ is diffeomorphic to an $S^2$ bundle over
$S^1$ or $M$ has a non-trivial connected sum decomposition. Thus, we
can establish results for all $3$-manifolds if we can establish them
for $3$-manifolds with $\pi_2=0$. Secondly, the reason for giving
this argument is that it is pleasing to see Ricci flow with surgery
implementing the connected sum decomposition required for
geometrization of $3$-manifolds. Also, this argument is a simpler
version of the one that we use to deal with components with
non-trivial $\pi_3$. Lastly, these results on Ricci flow do not use
the sphere theorem so that establishing the cutting into pieces with
trivial $\pi_2$ allows us to give a different proof of this result
(though admittedly one using much deeper ideas).



Let us now fix $T<\infty$ such that for all  $t\ge T$ all the
time-slices $M_t$ have trivial $\pi_2$. There is a simple
topological consequence of this and our assumption on the initial
manifold. If $M$ is a compact $3$-manifold whose fundamental group
is either a non-trivial free product or an infinite cyclic group,
then $M$ admits a homotopically non-trivial embedded $2$-sphere.
Since we began with a manifold $M_0$ whose fundamental group is a
free product of finite groups and infinite cyclic groups, it follows
that for $t\ge T$ every component of $M_t$ has finite fundamental
group.  Fix $t\ge T$. Then  each component of $M_t$ has a finite
cover that is simply connected, and thus, by an elementary argument
in algebraic topology, each component of $M_t$ has non-trivial
$\pi_3$. The second step in the finite-time extinction argument is
to use a non-trivial element in this group analogously to the way we
used homotopically non-trivial $2$-spheres to show that eventually
the manifolds have trivial $\pi_2$.

 There are two
approaches to this second step: the first is due to Perelman in
\cite{Perelman3} and the other due to Colding-Minicozzi in
\cite{ColdingMinicozzi}. In their approach Colding-Minicozzi
associate to a non-trivial element in $\pi_3(M)$ a non-trivial
element in $\pi_1(\mathrm{Maps}(S^2,M))$. This element is represented
by  a one-parameter family of $2$-spheres (starting and ending at
the constant map) representing a non-trivial element $\xi\in
\pi_3(M_0)$. They define the {\em width} of this homotopy class by
$W(\xi,t)$ by associating to each representative the maximal energy
of the $2$-spheres in the family and then minimizing over all
representatives of the homotopy class. Using results of Jost
\cite{Jost}, they show that this function satisfies the same forward
difference inequality that $W_2$ satisfies (and has the same
continuity property under Ricci flow and the same semi-continuity
under surgery). Since $W(\xi,t)$ is always $\ge 0$ if it is defined,
this forward difference quotient inequality implies that the
manifolds $M_t$ must eventually become empty.

While this approach seemed completely natural to us, and while we
believe that it works, we found the technical details
daunting\footnote{Colding and Minicozzi tell us they plan to give an
expanded version of their argument with a more detailed proof.}
(because one is forced to consider index-one critical points of the
energy functional rather than minima). For this reason we chose to
follow Perelman's approach. He represents a non-trivial element in
$\pi_3(M)$ as a non-trivial element in $\xi\in \pi_2(\Lambda M,*)$
where $\Lambda M$ is the free loop space of $M$. He then associates
to a family $\Gamma\colon S^2\to \Lambda M$ of homotopically trivial
loops an invariant $W(\Gamma)$ which is the maximum of the areas of
minimal spanning disks for the loops $\Gamma(c)$ as $c$ ranges over
$S^2$. The invariant of a non-trivial homotopy class $\xi$ is then
the infimum over all representatives $\Gamma$ for $\xi$ of
$W(\Gamma)$. As before, this function is continuous under Ricci flow
and is lower semi-continuous under surgery (unless the surgery
removes the component in question). It also satisfies a forward
difference quotient
$$\frac{dW(\xi)}{dt}\le -2\pi+\frac{3}{4t+1}W(\xi).$$
The reason for the term $-2\pi$ instead of $-4\pi$ which occurs in
the other cases is that we are working with minimal $2$-disks
instead of minimal $2$-spheres. Once this forward difference
quotient estimate (and the continuity) have been established the
argument finishes in the same way as the other cases: a function $W$
with the properties we have just established cannot be non-negative
for all positive time. This means the component in question, and
indeed all components at later time derived from it, must disappear
in finite time. Hence, under the hypothesis on the fundamental group
in Theorem~\ref{finiteext} the entire manifold must disappear at
finite time.

Because this approach uses only minima for the energy or area
functional, one does not have to deal with higher index critical
points. But one is forced to face other difficulties though --
namely boundary issues. Here, one must prescribe the deformation of
the family of boundary curves before computing the forward
difference quotient of the energy. The obvious choice is the
curve-shrinking flow (see \cite{AG}). Unfortunately, this flow can
only be defined when the curve in question is immersed and even in
this case the curve-shrinking flow can develop singularities even if
the Ricci flow does not. Following Perelman, or indeed \cite{AG},
one uses the device of taking the product with a small circle and
using loops, called {\em ramps}, that go around that
circle once. In this context the curve-shrinking flow remains
regular as long as the Ricci flow does. One then projects this flow
to a flow of families of $2$-spheres in the free loop space of the
time-slices of the original Ricci flow. Taking the length of the
circle sufficiently small yields the boundary deformation needed to
establish the forward difference quotient result. This requires a
compactness result which holds under local total curvature bounds.
This compactness result holds off of a set of time of small total
measure, which is sufficient for the argument. At the very end of
the argument we need an elementary but complicated result on annuli,
which we could not find in the literature. For more details on these
points see Chapter~\ref{energy}.


\section{Acknowledgements}

In sorting out Perelman's arguments we have had the aid of many
people. First of all, during his tour of the United States in the
Spring of 2003, Perelman spent enormous amounts of time explaining
his ideas and techniques both in public lectures and in private
sessions to us and to many others. Since his return to Russia, he
has also freely responded to our questions by e-mail. Richard
Hamilton has also given unstintingly of his time, explaining to us
his results and his ideas about Perelman's results. We have
benefitted tremendously from the work of Bruce Kleiner and John
Lott. They produced a lengthy set of notes filling in the details in
Perelman's arguments \cite{KleinerLott}. We have referred to those
notes countless times as we came to grips with Perelman's ideas. In
late August and early September of 2004, Kleiner, Lott and the two
of us participated in  a workshop at Princeton University, supported
by the Clay Math Institute, going through Perelman's second paper
(the one on Ricci flow with surgery) in detail. This workshop played
a significant role in convincing us that Perelman's arguments were
complete and correct and also in convincing us to write this book.
We thank all the participants of this workshop and especially
Guo-Feng Wei, Peng Lu, Yu Ding, and X.-C. Rong  who, together with
Kleiner and Lott, made significant contributions to the workshop.
Before, during, and after this workshop, we have benefitted from
private conversations too numerous to count with Bruce Kleiner and
John Lott.

Several of the analytic points were worked out in detail by others,
and we have freely adapted their work. Rugang Ye wrote a set of
notes proving the Lipschitz properties of the length function and
proving the fact that the cut locus for the reduced length function
is of measure zero. Closely related to this, Ye gave detailed
arguments proving the requisite limiting results for the length
function required to establish the existence of a gradient shrinking
soliton. None of these points was directly addressed by Perelman,
and it seemed to us that they needed careful explanation. Also, the
proof of the uniqueness of the standard solution was established
jointly by the second author and Peng Lu, and the proof of the
refined version of Shi's theorem where one has control on a certain
number of derivatives at time zero was also shown to us by Peng Lu.
We also benefitted from Ben Chow's expertise and vast knowledge of
the theory of Ricci flow especially during the ClayMath/MSRI Summer
School on Ricci Flow in 2005. We had several very helpful
conversations with Tom Mrowka and also with Robert Bryant,
especially about annuli of small area.


The second author gave courses at Princeton University and ran
seminars on this material at MIT and Princeton. Natasa Sesum and
Xiao-Dong Wang, see \cite{STW}, wrote notes for the seminars at MIT,
and Edward Fan and Alex Subotic took notes for the seminars and
courses at Princeton, and they produced preliminary manuscripts. We
have borrowed freely from these manuscripts, and it is a pleasure to
thank each of them for their efforts in the early stages of this
project.

The authors thank all the referees for their time and effort which
helped us to improve the presentation. In particular, we wish to
thank Cliff Taubes who carefully read the entire manuscript and gave
us enumerable comments and suggestions for clarifications and
improvements. We also thank Terry Tao for much helpful feedback
which improved the exposition. We also thank Colin Rourke for his
comments on the finite-time extinction argument.


We thank Lori Lejeune for drawing the figures and John Etnyre for technical
help in inserting these figures into the manuscript.

Lastly, during this work we were both generously supported by the Clay
Mathematical Institute, and it is a pleasure to thank the Clay Mathematics
Institute for its support and to thank its staff, especially Vida Salahi, for
their help in preparing this manuscript. We also thank the National Science Foundation
for their continued support and the second author thanks the Jim Simons
for his support during the period that he was a faculty member at MIT.

\section{List of related papers}

For the readers' convenience we gather here references to all the
closely related articles.

First and foremost are Perelman's three preprints, \cite{Perelman1},
\cite{Perelman2}, and \cite{Perelman3}. The first of these
introduces the main techniques in the case of Ricci flow, the second
discusses the extension of these techniques to Ricci flow with
surgery, and the last gives the short-cut to the Poincar\'e
Conjecture and the $3$-dimensional spherical space-form conjecture,
avoiding the study of the limits as time goes to infinity and
collapsing space arguments. There are the detailed notes by Bruce
Kleiner and John Lott, \cite{KleinerLott}, which greatly expand and
clarify Perelman's arguments from the first two preprints. There is
also a note on Perelman's second paper by Yu Ding  \cite{yuding}.
There is the article by Colding-Minicozzi \cite{ColdingMinicozzi},
which gives their alternate approach to the material in Perelman's
third preprint. Collapsing space arguments needed for the full
geometrization conjecture are discussed in Shioya-Yamaguchi
\cite{SY}. Lastly, after we had submitted a preliminary version of
this manuscript for refereeing, H.-D. Cao and X.-P. Zhu published an
article on the Poincar\'e Conjecture and Thurston's Geometrization
Conjecture; see \cite{CaoZhu}.
